\documentclass[10pt,journal,compsoc]{IEEEtran}

\usepackage{amsmath}
\usepackage{amsfonts}
\usepackage{amssymb}
\usepackage{booktabs}
\usepackage{hyperref}
\usepackage{graphicx}
\usepackage{microtype}
\usepackage{cite}

\hypersetup{
    colorlinks=true,
    linkcolor=blue,
    filecolor=magenta,      
    urlcolor=cyan,
    pdftitle={Operationalizing Semantic Data Engineering for Algorithmic Trade Execution},
    pdfpagemode=FullScreen,
}

\begin{document}

\title{Operationalizing Semantic Data Engineering for Algorithmic Trade Execution}

\author{Yufeng Chen, Marc Delgado, Manuel Rucker\\
\small{Department of Advanced Databases (BDA-GIA), Universitat Polit\`ecnica de Catalunya (UPC)}\\
\small{Date: May 23, 2026}}

\markboth{Semantic Data Management for AI, Project 2, May 2026}{}

\IEEEtitleabstractindextext{%
\begin{abstract}
Traditional quantitative finance pipelines rely on isolated tabular datasets and short-horizon binary classifiers, which fail to capture the relational structure of corporate markets and suffer from high transaction costs and silent failure modes when deployed. This paper presents an end-to-end operational semantic trading infrastructure. We construct a multi-layered financial and macroeconomic Knowledge Graph (KG) in our Exploitation Zone, linking corporate structures, sectors, and World Bank macroeconomic indicators via RDFLib. On top of this data fabric we deploy two analytical pipelines: a multi-graph SPARQL querying engine and a graph-neural-network pipeline in which RotatE structural embeddings are PCA-compressed and merged with technical and macroeconomic indicators to \emph{regress the 30-day cross-sectional rank} of each stock. The pivot to a relative-ordering target unlocks the structural KG signal: across an expanding-window walk-forward CV with $K{=}5$ folds and a 30-day embargo, the best tabular-only ensemble reaches a fold-mean Spearman information coefficient (IC) of only $+0.048$ (SE $0.023$), KG-only features climb to $+0.098$ (SE $0.029$), and the combined diverse soft-vote ensemble lands at $+0.105$ (SE $0.019$, $5/5$ folds positive) with a long-short decile spread of $+7.27\%$ per month. The semantic ranker drives an institutional-grade operational layer: a 2-state Gaussian Hidden Markov Model gates regime exposure, a state-space Kalman Beta Filter sizes systematic risk, an inverse-volatility rule sizes a concentrated top-K=10 book, and a differential portfolio optimiser executes only the weight deltas at a strict 10 basis points transaction-cost charge. On the strictly-post-training out-of-sample window the deployed book returns $+47.28\%$ versus $+22.18\%$ Buy~\&~Hold (Information Ratio $+4.37$). We further promote real-time news into a leak-controlled point-in-time sentiment feature, which lifts the combined IC to $+0.114$ and the out-of-sample return to $+57.49\%$, reported with explicit caveats. The live agent is hardened with three production controls absent from typical academic prototypes: a reproducibility hash that stamps every order with its model fingerprint and feature snapshot, a peak-to-trough drawdown circuit breaker, and a three-bucket HMM/ML/Kalman performance attribution log.
\end{abstract}

\begin{IEEEkeywords}
Semantic Data Engineering, Knowledge Graphs, Relational Graph Embeddings, Hidden Markov Models, Kalman Filters, Portfolio Rebalancing, Algorithmic Trading.
\end{IEEEkeywords}}

\maketitle
\IEEEdisplaynontitleabstractindextext
\IEEEpeerreviewmaketitle

\section{Introduction}
\IEEEPARstart{I}{n} quantitative asset management, predicting asset price movements is only half the battle. A highly accurate machine learning model is of zero economic value if it cannot be operationalized within a robust data engineering infrastructure. Data pipelines in the real world must ingest heterogeneous data, homogenize it semantically, manage systemic risk dynamically, and execute trades while minimizing transaction costs, spreads, and market slippage.

The academic objective of this project is to extend a classical tabular machine learning pipeline (developed in Phase 1) into an end-to-end semantic data science pipeline (Phase 2). P1 established a structured data lake (Landing, Formatted, and Trusted zones) to clean historical NASDAQ tick prices and train baseline Multi-Layer Perceptrons (MLP) and Random Forests. P2 transforms this architecture by embedding a **Knowledge Graph (KG)** into the Exploitation Zone. 

However, we recognized that simply feeding static relational embeddings into a classifier is insufficient to add real-world value. Therefore, we went far beyond the standard guidelines. We developed a complete, production-grade **Algorithmic Trading Bot CLI runner** that operates under a strict **10 basis points (10 bps)** transaction fee model. It dynamically adjusts market exposures using a **Gaussian Hidden Markov Model (HMM)** and scales short exposures using a time-varying **Kalman Beta Filter**—creating an institutional-grade, closed-loop trading infrastructure.

\section{System Architecture and Design Rationale}
\label{sec:arch}
Before detailing each component, we give the end-to-end map and---because every block was a deliberate choice over named alternatives---the reasoning behind it. The system is a single directed flow from raw data to broker orders:
\begin{center}
\small
\textbf{Landing} $\to$ \textbf{Formatted} $\to$ \textbf{Trusted} (Denial Constraints) $\to$ \textbf{Exploitation} (Knowledge Graph) $\to$ \{SPARQL analytics; RotatE embeddings\} $\to$ \textbf{feature fusion} $\to$ \textbf{diverse-ensemble rank model} $\to$ \textbf{operational overlay} (HMM regime $\cdot$ Kalman $\beta$ $\cdot$ inverse-vol sizing $\cdot$ differential rebalancer) $\to$ \textbf{Alpaca orders}.
\end{center}

Two principles drove the design. \emph{First, separation of concerns}: the supervised model answers only ``which stocks rank highest over the next 30 days,'' and every other economic judgement---when to be in the market, how much systematic risk to carry, how to split capital, how to trade cheaply---is handled by an explicit, individually-testable layer rather than being implicitly baked into one black-box predictor. \emph{Second, honesty under deployment}: each choice is validated under a leak-aware walk-forward protocol, and components that cannot be validated without look-ahead (e.g.\ live news as a tilt, before the PIT feature existed) are quarantined from the backtests. Table~\ref{tab:design} summarises the decisions; the rest of the paper substantiates each one.

\begin{table*}[tbp]
\centering
\caption{Architectural decisions and their rationale. Each layer was selected over concrete alternatives for the reason in the final column; the section that substantiates it is cited inline.}
\label{tab:design}
\small
\renewcommand{\arraystretch}{1.25}
\begin{tabular}{@{}p{2.5cm}p{3.4cm}p{10.3cm}@{}}
\toprule
\textbf{Layer / Decision} & \textbf{Choice made} & \textbf{Why this choice (and the alternative we rejected)} \\
\midrule
Data fabric & RDF Knowledge Graph in the Exploitation Zone (\S\ref{sec:arch}ff.) & Markets are inherently \emph{relational} (sector, industry, supply-chain, sovereign exposure). Flat SQL tables cannot express multi-hop risk paths; the KG lets a single SPARQL query traverse company\,$\to$\,country\,$\to$\,border\,$\to$\,geopolitical-tension, and gives the ML pipeline a structural substrate to embed. \\
Structural embedding & RotatE (complex-space rotations) & RotatE represents relations as unit-modulus rotations, so it can model \emph{symmetric}, \emph{inverse}, and \emph{compositional} relations. We rejected TransE (cannot model symmetry/1-to-N) and plain DistMult (cannot model anti-symmetry); both lose corporate-taxonomy structure RotatE preserves. \\
Dimensionality & PCA $256\!\to\!16$ ($\approx$79\% variance) & The raw 256-dim embedding would swamp the 41 tabular features and overfit on noisy financial data. PCA balances the two signal sources; we verified empirically that the \emph{uncompressed} fusion underperformed tabular-only, while the compressed fusion strictly improved it. \\
Prediction target & Cross-sectional 30-day \emph{rank} (PERCENT\_RANK) & The original 7-day binary direction is dominated by microstructure noise the (near-static) KG cannot help with. A relative 30-day rank is drift-invariant, matches how a long-short book is actually built, and is the target on which the KG signal finally pays off (IC $0.04\!\to\!0.11$). \\
Validation & Expanding walk-forward CV, $K{=}5$, 30-day embargo & A single 80/20 split cannot separate IC $0.07$ from $0.04$ on 3 months of test data, and ignores temporal leakage. The embargo equals the label horizon so no train row's forward window overlaps a test row. \\
Predictor & Diverse soft-vote ensemble (decorrelated trees + MLP) & No single learner dominated across folds; a correlation-pruned ensemble (members kept only if pairwise OOF corr $<0.85$) is more stable than any one model. We kept a Ridge and MLP baseline so the gain is argued against a fair non-linear floor, not a missing one. Rank-aware losses (lambdarank etc.) underperformed MSE and were dropped. \\
Regime gate & 2-state Gaussian HMM on S\&P 500 & Buying through a structural bear market is the dominant drawdown source. An unsupervised HMM infers latent bull/bear states from index returns and disables shorts in bull regimes---a market-wide judgement that does not belong in a per-stock ranker. \\
Systematic risk & State-space Kalman $\beta$ filter & Rolling-window $\beta$ lags regime shifts; a Kalman filter tracks time-varying $\beta$ recursively and reacts to a shock on the same bar, letting us shrink size on high-$\beta$ names instantly. \\
Position sizing & Inverse-volatility, capped at 25\% (\S\ref{sec:bt}) & On a concentrated book the dominant risk is \emph{idiosyncratic}, which HMM/Kalman (systematic-risk layers) do not address. $1/\sigma$ sizing with a per-name cap provably stops one low-vol name dominating the book. \\
Portfolio breadth & Full-universe scoring, hold top-K=10 & Scoring all $\approx$1{,}890 names preserves the cross-section the model was trained on; trading only the top 10 concentrates conviction. A breadth experiment (K=20 + sector cap) was \emph{worse} because the model's skill is sector-conditional---documented in \S\ref{sec:bt}, not hidden. \\
Execution & Differential rebalancer, 10 bps friction & Liquidate-and-rebuy burns spreads; trading only weight \emph{deltas} cut turnover $\sim$58\%. Sells execute before buys to free margin. \\
News signal & PIT news features in the model (was: live tilt) & News is genuinely predictive but needs point-in-time handling; we first shipped it as a live-only tilt (un-backtestable), then built strict as-of features and validated them (\S\ref{sec:bt}). \\
\bottomrule
\end{tabular}
\end{table*}

\section{The Semantic Data Engineering Pipeline}
The system's backbone is a structured, semantic DataOps pipeline. The data flows sequentially across four defined zones:

\subsection{Landing, Formatted, and Trusted Zones}
Raw data is ingested via Python collectors from Yahoo Finance, Nasdaq, and currency APIs, storing raw CSV logs in the \textbf{Landing Zone}. The \textbf{Formatted Zone} uses PySpark and DuckDB to syntactically homogenize these files into relational tables. The \textbf{Trusted Zone} applies strict quality rules expressed as Denial Constraints (e.g., ensuring closing prices are positive and volumes are non-negative). These constraints are both \emph{defined and executed} as Spark SQL views — the cleaned rows are collected from the Spark engine via \texttt{toPandas()} purely for persistence — and written, schema-parity-preserved against the Formatted Zone, to \texttt{TrustedZone.duckdb}.

\subsection{Exploitation Zone: The Knowledge Graph Layer}
Rather than relying on flat SQL tables, the Phase 2 \textbf{Exploitation Zone} models financial markets as linked data. We construct two distinct Knowledge Graphs in Terse RDF Triple Language (Turtle) formats:
\begin{enumerate}
    \item \textbf{Financial Knowledge Graph (\texttt{financial\_knowledge\_graph.ttl})}: Models the corporate taxonomy. Nodes include \texttt{Company}, \texttt{Sector}, \texttt{Industry}, \texttt{SizeClass}, \texttt{VolatilityClass}, and \texttt{Observation}. Edges model semantic relationships like \texttt{operatesInSector}, \texttt{belongsToIndustry}, and \texttt{headquarteredIn}. Each daily \texttt{Observation} carries close prices and trading volumes; the downstream tabular layer materialises both \texttt{target7dUp} (short-horizon binary direction, kept for backwards compatibility) and \texttt{target30dRank}, the cross-sectional \texttt{PERCENT\_RANK} of the forward 30-day return within each trading date. The country field for every ticker is resolved \emph{within the Exploitation Zone} as a cross-source reconciliation (sp500 country $\rightarrow$ Forbes name-match $\rightarrow$ default US) so the KG never receives NULL HQ countries; this enrichment is deliberately kept out of the Trusted Zone, which preserves strict schema parity with the Formatted Zone (per-source cleaned tables with Denial Constraints applied, no cross-source joins).
    \item \textbf{Macroeconomic \& Geopolitical Graph (\texttt{macroeconomic\_graph.ttl})}: Integrates broad-market context. It dynamically queries the RESTCountries API and the World Bank API to fetch population, geographical borders (\texttt{sharesBorderWith}), GDP, real GDP growth, and inflation rate literals, while encoding country-level geopolitical tension triples (\texttt{hasTensionWith}).
\end{enumerate}

By using a unified namespace, the financial and macroeconomic graphs are seamlessly linked at runtime, resolving ticker-level observations straight to sovereign-level indicators.

\section{SPARQL Pattern-Matching Analysis (Pipeline 1)}
To demonstrate the power of semantic data management, the first analytical pipeline leverages SPARQL to perform complex pattern matching. By modeling corporate and sovereign relationships as graphs, we query multidimensional risk paths that are extremely computationally expensive to write in SQL. 

We deployed ten analytical SPARQL queries. Most notably, \textbf{Q2}, \textbf{Q4}, and \textbf{Q7} demonstrate \textbf{cross-graph querying} by linking the financial and macroeconomic graphs. For example, \textbf{Q7} identifies companies headquartered in countries sharing a direct physical border with a nation experiencing geopolitical tensions:
\begin{equation}
\text{HQ} \xrightarrow{\text{headquarteredIn}} \text{Country}_1 \xrightarrow{\text{sharesBorderWith}} \text{Country}_2 \xrightarrow{\text{hasTensionWith}} \text{Rival}
\end{equation}
This allows risk compliance officers to audit indirect supply-chain and border proximity risk instantly, showcasing how semantic relationships add direct corporate value.

\section{Relational Graph Embeddings and Machine Learning (Pipeline 2)}
The second analytical pipeline translates our graph topology into mathematical feature vectors to train machine learning models.

\subsection{RotatE Structural Embedding Training}
Instead of basic TransE models which fail to capture symmetric or compositional relations, we implement \textbf{RotatE}~\cite{sun2019rotate}. RotatE models relations as rotations in a complex vector space. Given a triple $(h, r, t)$ representing head, relation, and tail, the entity embeddings $\mathbf{h}, \mathbf{t} \in \mathbb{C}^d$ are linked via the relation rotation $\mathbf{r} \in \mathbb{C}^d$ with $|\mathbf{r}_i| = 1$. The distance score function is defined as:
\begin{equation}
f_r(\mathbf{h}, \mathbf{t}) = \|\mathbf{h} \circ \mathbf{r} - \mathbf{t}\|_2
\end{equation}
where $\circ$ represents the Hadamard element-wise product. We fit a 128-dimensional complex space ($256$ real dimensions) using PyTorch, negative sampling, and a sigmoid-log loss.

\subsection{PCA Dimensionality Reduction \& Feature Merging}
To balance the high-dimensional structural embeddings against the small set of tabular indicators, we project the RotatE vectors into a lower-dimensional subspace using Principal Component Analysis:
\begin{equation}
\max_{\mathbf{W}} \text{Tr}(\mathbf{W}^T \mathbf{\Sigma} \mathbf{W}) \quad \text{s.t.} \quad \mathbf{W}^T \mathbf{W} = \mathbf{I}
\end{equation}
where $\mathbf{\Sigma}$ is the covariance matrix of the RotatE entity vectors. Projecting to 16 principal components retains $\approx 79\%$ of the structural variance while preventing the 256 raw embedding dimensions from swamping the 41-feature tabular signal. We empirically validated this trade-off: the uncompressed combined feature set yielded \emph{worse} cross-validated performance than tabular-only, whereas the PCA-compressed combined set yields strict improvements.

These 16 components are merged with 41 strictly past-only daily technical and macroeconomic indicators (1/5/10/20/50-day cumulative returns, vol-adjusted return, 5/10/20-day rolling volatility, RSI-14, MACD, Bollinger-band width, lagged returns, GDP, GDP growth, inflation, trade openness, policy rate) to form the consolidated feature matrix.

\subsection{Cross-Sectional 30-Day Rank Prediction}
Our earlier 7-day binary classification target proved a poor fit for graph-structural features: short-horizon direction is dominated by microstructure (momentum, volume shocks, vol regime), and the KG mostly encodes near-static corporate taxonomy. Following this empirical diagnosis, we pivoted the supervised target to the \textbf{cross-sectional 30-day rank}:
\begin{equation}
y_{i, t} = \text{PERCENT\_RANK}\!\left( r^{(30)}_{j, t} \right)_{j \in \mathcal{U}_t}, \qquad y_{i, t} \in [0, 1]
\end{equation}
where $r^{(30)}_{j, t}$ is the 30-day forward return for ticker $j$ on date $t$ and $\mathcal{U}_t$ is the universe of tickers trading on $t$. This is the natural target for portfolio construction: only the relative ordering matters, and the rank label is invariant to market-wide drift.

\subsection{Walk-Forward Cross-Validation with Embargo}
A single 80/20 temporal split on three months of test data cannot reliably distinguish an IC of $0.07$ from $0.04$. We therefore evaluate every (feature set, model) combination under an \textbf{expanding-window walk-forward CV} with $K=5$ folds and a \textbf{30-day embargo} between train and test slabs, matching the target horizon. Fold $k$ trains on all data up to the end of block $k$ and evaluates on block $k+1$. Per-fold metrics are aggregated as $\mu \pm \sigma$ to give an honest confidence interval. We report four portfolio-relevant metrics:
\begin{itemize}
    \item \textbf{Spearman IC} — rank correlation between predicted rank and realised rank.
    \item \textbf{Hit rate} — sign agreement between $(\hat{y}-0.5)$ and $(y-0.5)$.
    \item \textbf{Decile spread} — mean realised 30-day return of the predicted top decile minus the predicted bottom decile, the long-short P\&L an idealised portfolio would capture.
    \item \textbf{RMSE} — calibration sanity check on the $[0,1]$ rank output.
\end{itemize}

\subsection{Multi-Model Bake-Off}
The bake-off evaluates seven base learners — a regularised linear baseline (\texttt{RidgeCV}), a small two-layer MLP regressor described below, and four tree-based regressors (RandomForest, CatBoost, XGBoost, LightGBM) — alongside three rank-aware variants of the gradient-boosted trees (LightGBM \texttt{lambdarank}, XGBoost \texttt{rank:pairwise}, CatBoost \texttt{YetiRankPairwise}) and three ensembles: a \textbf{StackingRegressor} with a Ridge meta-learner, a rank-averaged \textbf{SoftVote} over the four tree regressors, and \textbf{SoftVote\_Diverse}, a greedy pruning of the SoftVote pool that keeps only members whose pairwise out-of-fold Pearson correlation stays below $0.85$. The MLP is a $64$-$64$ network with BatchNorm, dropout $0.2$, AdamW optimisation, and a temporal $90/10$ early-stopping split on each fold's training slab. The tree regressors use temporal $90/10$ early stopping inside each fold; tree depth and learning rate are kept conservative (\texttt{depth=4}, \texttt{lr=0.02}) to control variance on noisy financial data. Features are standardised cross-sectionally within each date (per-date z-score) rather than with a global \texttt{StandardScaler}, mirroring the standard quant-finance neutralisation; the live trading bot and the backtests apply the same per-date transform at inference.

\subsection{Cross-Sectional Leakage Analysis}
Several engineered features rank a ticker against its same-date peers (\texttt{PERCENT\_RANK() OVER (PARTITION BY Date ORDER BY $\cdot$)}). These features are safe under live deployment because: (a) each ranked input is itself strictly past-only (no \texttt{LEAD} window enters their computation); (b) every ticker's end-of-day close on date $t$ is published before the next session opens, so the cross-section is observable in real time; and (c) the target's forward 30-day window does not enter any feature computation. The 12-month master dataset spans 404{,}948 observations across 1{,}890 NASDAQ tickers from March 2025 to February 2026, after a quality filter that excludes penny stocks (close < \$1) and winsorises forward returns at $\pm 200\%$.

\subsection{Walk-Forward Bake-Off Results}
Table~\ref{tab:walk_forward_full} reports the full $K=5$ walk-forward bake-off (mean and SE across folds, plus the count of folds in which each metric was positive). On the tabular-only configuration the predictive signal is modest (best model XGBoost at $+0.0482$ IC, SE $0.0229$). Structural KGE features alone carry substantially more signal (SoftVote\_Diverse reaches $+0.0983$ IC, SE $0.0285$). Fusing tabular and embedding features delivers the best combination of signal and stability: SoftVote\_Diverse achieves $+0.1053$ IC (SE $0.0193$, 5/5 folds positive) and Stack maximises the long/short decile spread at $+7.66\%$ (SE $1.58\%$, 5/5 folds positive). Cross-sectional Z standardisation (Part 4 P3) and the small MLP regressor (Part 4 P4) participate in the bake-off so the paper can argue from a fair non-linear baseline, not a missing one.

\begin{table*}[tbp]
\centering
\caption{Walk-Forward Bake-Off summary. Mean and standard error ($\text{SE}=\sigma/\sqrt{K}$) across $K=5$ expanding folds, with a 30-day embargo between train and test. The IC sign and DS sign columns report how many of the $K$ folds had a positive metric (a one-sided binomial sign test against $\text{H}_0=0.5$).}
\label{tab:walk_forward_full}
\small
\begin{tabular}{llcccc}
\toprule
\textbf{Feature Set} & \textbf{Model} & \textbf{Spearman IC mean (SE)} & \textbf{IC sign} & \textbf{Decile Spread mean (SE)} & \textbf{DS sign} \\
\midrule
tabular only           & SoftVote           & +0.0406 (0.0238)                 & 4/5 & +4.76\% (1.15\%)                 & 5/5 \\
                       & SoftVote\_Rank     & -0.0179 (0.0282)                 & 2/5 & +2.01\% (1.66\%)                 & 3/5 \\
                       & SoftVote\_Diverse  & +0.0430 (0.0225)                 & 4/5 & +4.54\% (1.23\%)                 & 5/5 \\
                       & LightGBM           & +0.0463 (0.0219)                 & 4/5 & +3.96\% (1.16\%)                 & 5/5 \\
                       & LightGBM\_Rank     & -0.0678 (0.0513)                 & 1/5 & +0.17\% (3.47\%)                 & 3/5 \\
                       & XGBoost            & \textbf{+0.0482 (0.0229)}        & 4/5 & \textbf{+5.15\% (1.03\%)}        & 5/5 \\
                       & XGBoost\_Rank      & +0.0166 (0.0177)                 & 3/5 & +1.88\% (0.60\%)                 & 4/5 \\
                       & CatBoost           & +0.0080 (0.0280)                 & 3/5 & +3.23\% (1.20\%)                 & 5/5 \\
                       & CatBoost\_Rank     & +0.0181 (0.0175)                 & 3/5 & +2.53\% (1.21\%)                 & 4/5 \\
                       & RandomForest       & +0.0402 (0.0167)                 & 4/5 & +3.99\% (0.81\%)                 & 5/5 \\
                       & Stack              & +0.0452 (0.0246)                 & 4/5 & +5.10\% (1.26\%)                 & 5/5 \\
                       & MLP                & +0.0397 (0.0176)                 & 5/5 & +3.66\% (1.36\%)                 & 5/5 \\
                       & Ridge              & +0.0250 (0.0086)                 & 5/5 & +1.65\% (1.06\%)                 & 4/5 \\
\midrule
embedding only         & SoftVote           & +0.0943 (0.0255)                 & 5/5 & +5.91\% (2.29\%)                 & 4/5 \\
                       & SoftVote\_Rank     & -0.0081 (0.0540)                 & 2/5 & +4.48\% (3.90\%)                 & 4/5 \\
                       & SoftVote\_Diverse  & \textbf{+0.0983 (0.0285)}        & 5/5 & +6.31\% (2.20\%)                 & 5/5 \\
                       & LightGBM           & +0.0811 (0.0223)                 & 5/5 & +6.10\% (1.92\%)                 & 5/5 \\
                       & LightGBM\_Rank     & -0.0454 (0.0590)                 & 1/5 & +1.78\% (3.48\%)                 & 4/5 \\
                       & XGBoost            & +0.0948 (0.0279)                 & 5/5 & +5.91\% (1.95\%)                 & 5/5 \\
                       & XGBoost\_Rank      & -0.0220 (0.0505)                 & 2/5 & +3.51\% (4.01\%)                 & 4/5 \\
                       & CatBoost           & +0.0855 (0.0271)                 & 5/5 & +6.43\% (1.80\%)                 & 5/5 \\
                       & CatBoost\_Rank     & +0.0583 (0.0292)                 & 4/5 & \textbf{+6.47\% (3.36\%)}        & 3/5 \\
                       & RandomForest       & +0.0885 (0.0179)                 & 5/5 & +5.13\% (1.45\%)                 & 5/5 \\
                       & Stack              & +0.0934 (0.0270)                 & 5/5 & +5.96\% (1.94\%)                 & 5/5 \\
                       & MLP                & +0.0884 (0.0301)                 & 4/5 & +6.15\% (1.58\%)                 & 5/5 \\
                       & Ridge              & +0.0492 (0.0375)                 & 4/5 & +5.02\% (2.91\%)                 & 3/5 \\
\midrule
tabular + embedding    & SoftVote           & +0.1007 (0.0176)                 & 5/5 & +7.20\% (1.51\%)                 & 5/5 \\
                       & SoftVote\_Rank     & -0.0171 (0.0543)                 & 2/5 & +3.24\% (3.71\%)                 & 4/5 \\
                       & SoftVote\_Diverse  & \textbf{+0.1053 (0.0193)}        & 5/5 & +7.27\% (1.66\%)                 & 5/5 \\
                       & LightGBM           & +0.0893 (0.0175)                 & 5/5 & +7.48\% (1.44\%)                 & 5/5 \\
                       & LightGBM\_Rank     & -0.0559 (0.0627)                 & 1/5 & -1.22\% (3.53\%)                 & 2/5 \\
                       & XGBoost            & +0.1031 (0.0219)                 & 5/5 & +7.52\% (1.41\%)                 & 5/5 \\
                       & XGBoost\_Rank      & -0.0179 (0.0575)                 & 2/5 & +3.24\% (3.70\%)                 & 3/5 \\
                       & CatBoost           & +0.0676 (0.0128)                 & 5/5 & +6.34\% (0.96\%)                 & 5/5 \\
                       & CatBoost\_Rank     & +0.0379 (0.0187)                 & 3/5 & +4.08\% (2.47\%)                 & 3/5 \\
                       & RandomForest       & +0.0999 (0.0135)                 & 5/5 & +5.89\% (1.50\%)                 & 5/5 \\
                       & Stack              & +0.1010 (0.0230)                 & 5/5 & \textbf{+7.66\% (1.58\%)}        & 5/5 \\
                       & MLP                & +0.0893 (0.0213)                 & 5/5 & +6.15\% (1.37\%)                 & 5/5 \\
                       & Ridge              & +0.0625 (0.0310)                 & 4/5 & +4.92\% (2.74\%)                 & 3/5 \\
\bottomrule
\end{tabular}
\end{table*}

\subsection{Rigorous Pairwise Statistical Inference}
Holding the model constant, we compare each feature configuration against the others. For each model and each comparison we take the $K=5$ paired fold-level differences, run a two-sided binomial sign test against $\text{H}_0=0.5$, and bootstrap a 95\% CI on the mean of those five differences from 10{,}000 paired resamples. With $K=5$ the minimum attainable two-sided sign-test $p$ is $0.0625$ at $5/5$ — the bootstrap CI is the complementary range estimate. We flag a row as significant only when both the CI excludes zero \emph{and} the sign-test $p<0.20$ (Table~\ref{tab:pairwise_stats}).

On the IC, the Combined-minus-Tabular contrast clears that bar for 8 of the 13 models tested, with paired-mean IC improvements ranging from $+0.0429$ to $+0.0624$ and 95\% bootstrap CIs entirely above zero. The Ridge baseline (Part 4 P1) is materially weaker: its Combined-minus-Tabular IC gain is $+0.0375$ with a bootstrap CI of $[-0.0036, +0.0797]$ and sign-test $p=0.375$ — that CI crosses zero, so we cannot reject the null that adding KG embeddings does nothing for Ridge. Non-linear ensembles are doing the heavy lifting; linearity alone does not extract the cross-feature interaction the KG embeddings encode. Rank-aware losses (LightGBM lambdarank, XGBoost rank:pairwise, CatBoost YetiRankPairwise) underperformed their MSE counterparts in every feature configuration and are reported in Table~\ref{tab:walk_forward_full} for transparency rather than recommended for deployment.

\begin{table*}[tbp]
\centering
\caption{Paired pairwise comparisons across the three feature configurations, holding the model constant. The mean delta is the average of $K=5$ paired fold-level differences; the 95\% confidence interval is a 10{{,}}000-resample paired bootstrap on those five differences; the sign-test $p$ is a two-sided binomial test on the signs of the differences. Bold rows are those whose CI excludes zero \emph{and} whose sign-test $p<0.20$.}
\label{tab:pairwise_stats}
\small
\begin{tabular}{llcccc}
\toprule
\textbf{Metric} & \textbf{Comparison} & \textbf{Model} & \textbf{Mean delta} & \textbf{95\% bootstrap CI} & \textbf{Sign-test } $p$ \\
\midrule
Decile Spread  & Embedding $-$ Tabular    & SoftVote           & +1.14\%                & [-2.73\%, +5.90\%]               & 1.000      \\
               &                          & SoftVote\_Rank     & +2.47\%                & [-2.34\%, +6.67\%]               & 1.000      \\
               &                          & SoftVote\_Diverse  & +1.78\%                & [-1.93\%, +5.84\%]               & 0.375      \\
               &                          & LightGBM           & +2.14\%                & [-1.82\%, +6.90\%]               & 1.000      \\
               &                          & LightGBM\_Rank     & +1.61\%                & [-0.48\%, +3.21\%]               & 0.375      \\
               &                          & XGBoost            & +0.76\%                & [-2.55\%, +4.79\%]               & 1.000      \\
               &                          & XGBoost\_Rank      & +1.63\%                & [-4.00\%, +7.68\%]               & 1.000      \\
               &                          & CatBoost           & +3.19\%                & [+0.16\%, +7.68\%]               & 0.375      \\
               &                          & CatBoost\_Rank     & +3.94\%                & [+0.02\%, +7.82\%]               & 1.000      \\
               &                          & RandomForest       & +1.14\%                & [-0.79\%, +4.37\%]               & 1.000      \\
               &                          & Stack              & +0.85\%                & [-2.26\%, +4.98\%]               & 1.000      \\
               &                          & MLP                & +2.49\%                & [+0.24\%, +4.74\%]               & 1.000      \\
               &                          & Ridge              & +3.37\%                & [-0.33\%, +7.07\%]               & 1.000      \\
               & Combined $-$ Embedding   & SoftVote           & +1.29\%                & [-2.13\%, +4.46\%]               & 0.375      \\
               &                          & SoftVote\_Rank     & -1.25\%                & [-2.93\%, +0.40\%]               & 0.375      \\
               &                          & SoftVote\_Diverse  & +0.96\%                & [-1.99\%, +3.85\%]               & 0.375      \\
               &                          & LightGBM           & +1.38\%                & [-2.26\%, +4.28\%]               & 0.375      \\
               &                          & LightGBM\_Rank     & \textbf{-3.00\%}       & \textbf{[-3.93\%, -2.31\%]}      & \textbf{0.062} \\
               &                          & XGBoost            & +1.61\%                & [-0.75\%, +3.38\%]               & 0.375      \\
               &                          & XGBoost\_Rank      & -0.27\%                & [-1.79\%, +1.24\%]               & 1.000      \\
               &                          & CatBoost           & -0.08\%                & [-2.87\%, +2.08\%]               & 1.000      \\
               &                          & CatBoost\_Rank     & -2.39\%                & [-3.88\%, -0.63\%]               & 0.375      \\
               &                          & RandomForest       & +0.76\%                & [-1.85\%, +3.26\%]               & 1.000      \\
               &                          & Stack              & +1.70\%                & [-0.31\%, +3.67\%]               & 0.375      \\
               &                          & MLP                & +0.01\%                & [-2.14\%, +2.15\%]               & 1.000      \\
               &                          & Ridge              & -0.10\%                & [-0.59\%, +0.54\%]               & 1.000      \\
               & Combined $-$ Tabular     & SoftVote           & +2.44\%                & [+0.86\%, +3.87\%]               & 0.375      \\
               &                          & SoftVote\_Rank     & +1.23\%                & [-3.45\%, +5.45\%]               & 1.000      \\
               &                          & SoftVote\_Diverse  & \textbf{+2.73\%}       & \textbf{[+1.33\%, +4.13\%]}      & \textbf{0.062} \\
               &                          & LightGBM           & \textbf{+3.52\%}       & \textbf{[+1.19\%, +6.18\%]}      & \textbf{0.062} \\
               &                          & LightGBM\_Rank     & -1.39\%                & [-3.76\%, +0.54\%]               & 0.375      \\
               &                          & XGBoost            & +2.37\%                & [+0.58\%, +4.00\%]               & 0.375      \\
               &                          & XGBoost\_Rank      & +1.36\%                & [-4.06\%, +6.77\%]               & 1.000      \\
               &                          & CatBoost           & \textbf{+3.11\%}       & \textbf{[+1.70\%, +4.89\%]}      & \textbf{0.062} \\
               &                          & CatBoost\_Rank     & +1.55\%                & [-0.83\%, +3.94\%]               & 1.000      \\
               &                          & RandomForest       & +1.90\%                & [+0.21\%, +3.44\%]               & 0.375      \\
               &                          & Stack              & \textbf{+2.56\%}       & \textbf{[+0.80\%, +4.63\%]}      & \textbf{0.062} \\
               &                          & MLP                & \textbf{+2.50\%}       & \textbf{[+2.19\%, +2.92\%]}      & \textbf{0.062} \\
               &                          & Ridge              & +3.27\%                & [-0.28\%, +6.81\%]               & 1.000      \\
Spearman IC    & Embedding $-$ Tabular    & SoftVote           & +0.0537                & [+0.0132, +0.0976]               & 0.375      \\
               &                          & SoftVote\_Rank     & +0.0097                & [-0.0532, +0.0610]               & 0.375      \\
               &                          & SoftVote\_Diverse  & +0.0553                & [+0.0165, +0.0941]               & 0.375      \\
               &                          & LightGBM           & +0.0347                & [-0.0202, +0.0897]               & 1.000      \\
               &                          & LightGBM\_Rank     & +0.0224                & [+0.0014, +0.0435]               & 0.375      \\
               &                          & XGBoost            & +0.0465                & [-0.0000, +0.0953]               & 0.375      \\
               &                          & XGBoost\_Rank      & -0.0386                & [-0.1339, +0.0352]               & 1.000      \\
               &                          & CatBoost           & \textbf{+0.0775}       & \textbf{[+0.0251, +0.1299]}      & \textbf{0.062} \\
               &                          & CatBoost\_Rank     & +0.0402                & [-0.0063, +0.0838]               & 0.375      \\
               &                          & RandomForest       & \textbf{+0.0483}       & \textbf{[+0.0231, +0.0769]}      & \textbf{0.062} \\
               &                          & Stack              & +0.0481                & [+0.0012, +0.0952]               & 0.375      \\
               &                          & MLP                & +0.0487                & [+0.0094, +0.0883]               & 0.375      \\
               &                          & Ridge              & +0.0242                & [-0.0318, +0.0754]               & 0.375      \\
               & Combined $-$ Embedding   & SoftVote           & +0.0065                & [-0.0211, +0.0352]               & 1.000      \\
               &                          & SoftVote\_Rank     & -0.0090                & [-0.0243, +0.0099]               & 0.375      \\
               &                          & SoftVote\_Diverse  & +0.0071                & [-0.0187, +0.0375]               & 1.000      \\
               &                          & LightGBM           & +0.0082                & [-0.0173, +0.0368]               & 1.000      \\
               &                          & LightGBM\_Rank     & -0.0105                & [-0.0176, -0.0033]               & 0.375      \\
               &                          & XGBoost            & +0.0083                & [-0.0135, +0.0335]               & 1.000      \\
               &                          & XGBoost\_Rank      & +0.0040                & [-0.0275, +0.0389]               & 1.000      \\
               &                          & CatBoost           & -0.0179                & [-0.0849, +0.0356]               & 1.000      \\
               &                          & CatBoost\_Rank     & -0.0204                & [-0.0500, +0.0039]               & 0.375      \\
               &                          & RandomForest       & +0.0115                & [-0.0064, +0.0338]               & 1.000      \\
               &                          & Stack              & +0.0076                & [-0.0149, +0.0351]               & 1.000      \\
               &                          & MLP                & +0.0010                & [-0.0183, +0.0299]               & 0.375      \\
               &                          & Ridge              & +0.0133                & [+0.0018, +0.0275]               & 0.375      \\
               & Combined $-$ Tabular     & SoftVote           & \textbf{+0.0601}       & \textbf{[+0.0438, +0.0795]}      & \textbf{0.062} \\
               &                          & SoftVote\_Rank     & +0.0008                & [-0.0633, +0.0441]               & 0.375      \\
               &                          & SoftVote\_Diverse  & \textbf{+0.0624}       & \textbf{[+0.0475, +0.0782]}      & \textbf{0.062} \\
               &                          & LightGBM           & \textbf{+0.0429}       & \textbf{[+0.0135, +0.0724]}      & \textbf{0.062} \\
               &                          & LightGBM\_Rank     & +0.0120                & [-0.0162, +0.0382]               & 1.000      \\
               &                          & XGBoost            & \textbf{+0.0549}       & \textbf{[+0.0303, +0.0815]}      & \textbf{0.062} \\
               &                          & XGBoost\_Rank      & -0.0345                & [-0.1482, +0.0462]               & 1.000      \\
               &                          & CatBoost           & \textbf{+0.0596}       & \textbf{[+0.0202, +0.0990]}      & \textbf{0.062} \\
               &                          & CatBoost\_Rank     & +0.0198                & [-0.0104, +0.0464]               & 0.375      \\
               &                          & RandomForest       & \textbf{+0.0597}       & \textbf{[+0.0438, +0.0757]}      & \textbf{0.062} \\
               &                          & Stack              & \textbf{+0.0557}       & \textbf{[+0.0321, +0.0811]}      & \textbf{0.062} \\
               &                          & MLP                & \textbf{+0.0496}       & \textbf{[+0.0266, +0.0788]}      & \textbf{0.062} \\
               &                          & Ridge              & +0.0375                & [-0.0036, +0.0797]               & 0.375      \\
\bottomrule
\end{tabular}
\end{table*}

\subsection{OOF Diversity Diagnostic and Ensemble Construction}
To understand whether SoftVote is averaging genuinely independent signal or just adding similar trees together, we compute the pairwise Pearson correlation of out-of-fold predictions on the combined feature set across the $K=5$ folds. The gradient-boosted trees (CatBoost, LightGBM, XGBoost) correlate with each other in the range $r\in[0.725, 0.877]$ — strong but not redundant. Ridge correlates with the trees only at $r\in[0.372, 0.505]$, which is exactly the diversity profile that lets a simple rank-average ensemble extract additional variance reduction. The diverse-ensemble selector (Part 4 P5) pruned the SoftVote pool to members whose pairwise OOF Pearson correlation stays below 0.85, yielding SoftVote\_Diverse at $+0.1053$ IC (SE $0.0193$) and $+7.27\%$ decile spread (SE $1.66\%$). The deployed artefact persisted to \texttt{best\_model.pkl} is the SoftVote ensemble over the gradient-boosted regressors on the combined feature set, with the cross-sectional Z preprocessing flag set so that the live trading bot and the backtests apply the same per-date normalisation used at training time.

\section{Operational Quantitative Infrastructure (Going Beyond)}
The primary contribution of our work is the translation of these semantic classifiers into an automated, operational quantitative execution infrastructure. The infrastructure is composed of three state-of-the-art risk-management and trading layers:

\subsection{Gaussian Hidden Markov Model (HMM) Regime switching}
Markets fluctuate between distinct volatility cycles. A predictive model trained to buy stocks in bull markets will suffer massive drawdowns if it continues to buy during structural bear markets. To address this, we implement a rolling \textbf{2-state Gaussian Hidden Markov Model} on daily S\&P 500 log returns. 

The S\&P 500 returns are modeled as emission probability distributions generated by unobserved latent market states $S_t \in \{0, 1\}$, where State 0 is characterized by low volatility (Bull) and State 1 by high volatility (Bear). The parameters (transition probabilities $A$, emission means and variances $B$, and initial probabilities $\pi$) are fitted dynamically using the Baum-Welch (Expectation-Maximization) algorithm. The optimal state sequence is decoded using the Viterbi trellis path:
\begin{equation}
V_{t}(j) = \max_{i} \left[ V_{t-1}(i) a_{ij} \right] b_j(x_t)
\end{equation}
If the current decoded state is State 1 (Bear), the system enables shorting. If State 0 (Bull) is decoded, short positions are automatically disabled to prevent margin squeezes.

\subsection{State-Space Kalman Beta Filter Risk Scaling}
To protect active capital from highly volatile assets, the bot measures the time-varying systemic risk ($\beta$) of each stock relative to the market index using a \textbf{Kalman Filter}. The system state is defined by the regression parameters:
\begin{equation}
\mathbf{\theta}_t = [\alpha_t, \beta_t]^T
\end{equation}
The transition and measurement equations are:
\begin{align}
\mathbf{\theta}_t &= \mathbf{\theta}_{t-1} + \mathbf{w}_t, \quad \mathbf{w}_t \sim \mathcal{N}(0, \mathbf{Q}) \\
y_t &= \mathbf{H}_t \mathbf{\theta}_t + v_t, \quad v_t \sim \mathcal{N}(0, R)
\end{align}
where $\mathbf{H}_t = [1, x_{\text{SP500}, t}]$. The filter recursively updates the state vector using the Kalman Gain $\mathbf{K}_t$ and the prediction error covariance $\mathbf{P}_t$. In Bear states, the target short allocations are dynamically scaled down for high-beta stocks:
\begin{equation}
W_{\text{target}, i} = \frac{W_{\text{raw}, i}}{\max(|\beta_{\text{Kalman}, i}|, 0.5)}
\end{equation}
This caps size exposure on high-systemic-risk stocks, compressing drawdowns.

\subsection{Differential Portfolio Rebalancing Optimizer}
Naive rebalancing (selling the entire portfolio to cash and buying the new target weights) generates massive transaction fee friction. Our bot implements a \textbf{Differential Portfolio Rebalancing Optimizer}. It computes only the absolute weight delta changes:
\begin{equation}
\Delta W_i = W_{\text{target}, i} - W_{\text{actual}, i}
\end{equation}
This simple optimization \textbf{reduces rebalancing trade volume by over 58.32\%}, directly saving spreads and commissions. Furthermore, the bot executes SELL and short-covering orders first, freeing up brokerage account margin before submitting BUY orders, preventing margin rejections.

\subsection{Inverse-Volatility Position Sizing}
The concentrated top-$K{=}10$ deployment introduces a risk that the HMM regime gate and Kalman beta filter do not address. Those two layers manage \emph{systematic} risk (market-wide co-movement, captured by $\beta$), which dominates a broadly diversified book; but on a ten-name book the dominant risk is \emph{idiosyncratic} single-name variance, which diversifies as $1/N$ and is therefore large when $N$ is small. To target that risk without diluting the model's stock selection, we separate concerns: the ranking model decides \emph{which} names enter the book, and an inverse-volatility rule decides \emph{how much} of each:
\begin{equation}
w_i \;\propto\; \frac{1}{\sigma_i}, \qquad w_i \le w_{\max},\qquad \sum_i w_i = 1,
\end{equation}
where $\sigma_i$ is the trailing 20-day realised volatility and $w_{\max}{=}0.25$ caps any single position (enforced by freezing capped names and redistributing the residual budget among the rest). High-volatility names thus receive proportionally less capital, so a single name's blow-up contributes less to portfolio drawdown.

On the strictly-out-of-sample post-training window (Section~\ref{sec:bt}), the inverse-volatility top-$K{=}10$ book returns $+47.28\%$ at Sharpe $6.66$, maximum drawdown $-2.40\%$, and an Information Ratio of $+4.37$ against Buy~\&~Hold. We justify the $1/\sigma$ rule primarily on \emph{risk-management} grounds rather than a backtested return edge: the per-name cap $w_{\max}{=}0.25$ provably prevents any single ultra-low-volatility name from dominating the book, which is the failure mode a ten-name portfolio is most exposed to. We deliberately do \emph{not} claim a quantified inverse-vol-versus-equal-weight improvement here: an earlier comparison rested on an 18-month window that we subsequently found overlaps the model's training range ($2025\text{-}03\text{-}31 \to 2026\text{-}01\text{-}14$), so those deltas were contaminated and have been withdrawn. The binding constraint on a ten-name book is the concentration itself; $1/\sigma$ sizing trims the edges of that risk without eliminating it, and materially lower drawdown would require relaxing $K$ --- which, as the breadth experiment in Section~\ref{sec:bt} shows, trades away the very cross-sectional alpha the concentrated book is designed to capture.

\subsection{News-Sentiment as a Semantic Feature (v1 Live Tilt $\to$ v2 Backtested PIT Feature)}
As a forward-looking semantic enhancement, the live agent ingests real-time company news and tilts the model's ranking by sentiment. At decision time the agent queries the Alpaca News API for each candidate, scores each headline/summary (a finance lexicon by default, with optional FinBERT), and --- staying within the project's knowledge-graph paradigm --- assembles an \emph{ephemeral} RDF \texttt{NewsEvent} graph in which each article node links to its company node carrying \texttt{sentimentScore} and \texttt{publishedAt} literals. A SPARQL query then aggregates mean sentiment per ticker, and the agent re-ranks:
\begin{equation}
\text{score}_i = \text{pred\_rank}_i + \lambda\, s_i, \qquad s_i \in [-1, 1],\ \lambda = 0.10,
\end{equation}
where $s_i$ is the aggregated news sentiment. The model still decides the base ordering; recent news only nudges it.

This tilt was the \emph{first} (v1) form of the feature, deliberately live-only: a leak-free historical study needs a point-in-time (PIT) archive --- each article timestamped at its true public-release moment --- and we initially treated that as future work.

\subsubsection{From Tilt to Backtested Feature (v2)}
We subsequently \emph{realized} that future work. Alpaca's News API returns each article with its true \texttt{created\_at} timestamp and supports historical date-range queries, which is enough to construct PIT features under a strict as-of rule. We therefore built five news features per (ticker, date) --- 7-day mean sentiment, news volume, 30-day mean sentiment, 7-vs-30-day sentiment momentum, and 7-day dispersion --- where the feature for date $D$ aggregates only articles published \emph{strictly before} $D$ (no same-day or future leak), scored with the \emph{frozen} lexicon so the scorer never sees test-period text. A single canonical implementation feeds three consumers (the training-set builder, live inference, and the backtests), so the definitions cannot diverge. Coverage spans $1{,}518/1{,}890$ tickers ($\sim$82k articles), though it is skewed toward large caps.

Joined into the tabular block and re-evaluated under the \emph{identical} walk-forward CV (same embeddings, data, folds; only $+5$ columns), news lifts the Combined IC of the diverse ensemble from $+0.1053$ to $+0.1135$ --- positive on all eight model families and $4/5$ folds, though within the $\sim\!0.043$ fold-to-fold standard deviation (a $4/5$ sign test gives $p\!\approx\!0.19$): directionally consistent but not individually conclusive. The portfolio-level test on the clean post-training OOS window is more decisive: the news-trained full-universe top-K=10 book returns \textbf{$+57.49\%$ vs $+47.28\%$} for the news-free model (Buy \& Hold $+22.18\%$), with Information Ratio rising to $+4.52$ (from $+4.37$), at marginally higher volatility (Sharpe $6.55$ vs $6.66$, Max DD $-3.39\%$ vs $-2.40\%$). We adopt the PIT news features into the deployed model and compute them live at inference time; the post-hoc tilt is correspondingly disabled (\texttt{USE\_NEWS\_SENTIMENT=0}) to avoid double-counting. Honest caveats stand: PIT integrity is \emph{assumed} from Alpaca timestamps (no vendor guarantee against silent back-fill), coverage is large-cap-skewed, entity resolution uses current symbols, and the portfolio gain rests on a single 2-month / 9-rebalance window. \textbf{All other backtests in this paper use the news-free model and are unaffected.}

\subsection{Operational Hardening for Live Deployment}
A bot that survives backtests but cannot be audited or stopped under real money is a liability rather than an asset. Beyond the headline trading logic, our agent ships four production controls collected in a dedicated \texttt{operational.py} module so they are reusable across agents:

\textbf{Reproducibility hash.} Every rebalance computes a \texttt{decision\_id} as the SHA-256 prefix of \texttt{(model\_hash, feature\_snapshot\_hash, universe\_hash, hmm\_state, decision\_date)}. The exact feature matrix used at decision time is persisted to \texttt{./agent\_logs/decisions/<decision\_id>.parquet}, so any trade can be replayed offline with the same model file and the same inputs.

\textbf{Per-trade structured log.} Each submitted order writes a row to \texttt{./agent\_logs/trades.csv} containing the order action, target and delta weight, intended quantity and notional, reference price, ML signal, HMM state, HMM bull probability, Kalman beta, the decision id, and a dry-run flag. This is the equivalent of an institutional FIX-log for post-trade attribution and root-cause analysis.

\textbf{Drawdown circuit breaker.} A running peak-to-trough equity tracker persisted to \texttt{equity\_history.csv} computes the live drawdown against the rolling peak. When the breach threshold (default $5\%$) is exceeded, the bot creates a halt file \texttt{.agent\_halted} containing the offending equity and timestamp, and refuses to submit any further orders until the file is removed by a human operator. The reset is deliberately manual: bots that auto-resume after a drawdown tend to chase losses into a deeper one.

\textbf{Three-bucket performance attribution.} Between rebalances we decompose realised P\&L into three approximate buckets:
\begin{equation}
\Delta\text{PnL}_t = \underbrace{(\Delta\text{Exposure}_t) \cdot r_{\text{SPX}, t}}_{\text{regime}} + \underbrace{\sum_i \tilde{w}_i\, \tilde{r}_i}_{\text{ML rank}} + \underbrace{\sum_i \beta_i\, r_{\text{SPX}, t}\, w_i}_{\text{Kalman}} + \text{residual}
\end{equation}
where $\tilde{w}_i$ and $\tilde{r}_i$ are the de-meaned weights and returns within the universe. The decomposition is intentionally approximate under correlated drivers — what matters is that we can read off whether a profitable week was driven by the regime call, the ML ranking, or the Kalman beta sizing. A persistently large residual flags that something we don't measure (idiosyncratic news, fills, market impact) is driving the book.

\section{Rigorous Out-of-Sample Validation and Slippage Simulation}
\label{sec:bt}
To evaluate the operational validity of the system under realistic institutional constraints, we implemented a strict **10 basis points (10 bps)** transaction fee and slippage model on all rebalancing turnover.

\subsection{Drift Propagation Mathematical Consistency}
Between weekly rebalancing points, the portfolio's assets drift with price returns. The actual weight of asset $i$ at the end of week $t$ is mathematically propagated as:
\begin{equation}
W_{\text{actual}, i, t} = W_{\text{target}, i, t-1} \times \frac{1 + R_{i, t}}{1 + R_{\text{portfolio}, t}}
\end{equation}
where $R_{\text{portfolio}, t} = \sum_j W_{\text{target}, j, t-1} R_{j, t}$. Weekly rebalancing turnover is computed as:
\begin{equation}
\text{Turnover}_t = \sum_{i \in \text{stocks}} |W_{\text{target}, i, t} - W_{\text{actual}, i, t}|
\end{equation}
The 10 bps transaction fee is strictly deducted from portfolio equity before weekly returns are compounded:
\begin{equation}
\text{Equity}_{\text{post-rebalance}, t} = \text{Equity}_{\text{pre-rebalance}, t} \times (1.0 - \text{Turnover}_t \times \text{FEE\_RATE})
\end{equation}
where $\text{FEE\_RATE} = 0.0010$. To ensure mathematical fairness, the Buy \& Hold (B\&H) benchmark is charged a 10 bps entry fee on Day 1 and a 10 bps exit fee on the final liquidation day.

\subsection{Empirical Backtest Results (30-Day Rank Model)}
The walk-forward CV reported earlier (IC $+0.1053$, $5/5$ folds positive) is a training-time diagnostic; the question this section answers is whether that signal survives a portfolio-level backtest under realistic 10 bps friction. Two deployment modes are evaluated to expose the gap between the two:

\begin{itemize}\itemsep0pt
  \item \textbf{Curated 20-ticker basket}: the historical \texttt{HIGH\_ALPHA\_TICKERS} alphabetical small-cap list used by earlier versions of the agent.
  \item \textbf{Full-universe top-K=10 ($\geq$ top 5\%)}: every Friday the model scores all $\approx$1{,}890 tickers it can embed, the top-K=10 with predicted rank $\geq 0.95$ are inverse-volatility weighted (Section above), and the differential rebalancer turns that into the next set of orders. This mode is consistent with the cross-sectional Z preprocessing the model was trained against ($N=1{,}890$ at training, $N=1{,}890$ at inference).
\end{itemize}

\noindent \textbf{What counts as out-of-sample.} The deployed model was fit on feature-dates $2025\text{-}03\text{-}31 \to 2026\text{-}01\text{-}14$ (the underlying data ends $2026\text{-}02\text{-}13$; the last 30 days are trimmed because the target is a 30-day forward return). A backtest is therefore genuinely out-of-sample (no memorisation) only on windows \emph{outside} that interval --- before $2025\text{-}03\text{-}31$ or after $2026\text{-}01\text{-}14$. Earlier drafts of this report headlined 6-, 12- and 24-month horizons whose windows overlap the training range; those returns reflect the model scoring data it was trained on and have been \textbf{removed}. We report only the two genuinely-OOS windows below.

\begin{table}[htbp]
\caption{Genuinely out-of-sample full-universe top-K=10 performance under 10 bps friction (inverse-volatility sizing). Both windows exclude the training feature-dates ($2025\text{-}03\text{-}31 \to 2026\text{-}01\text{-}14$). IR is the annualised Information Ratio against Buy~\&~Hold. $^{\ddagger}$ The pre-training window additionally carries current-membership \emph{survivorship bias} and mild \emph{static-embedding look-ahead} (the KG structural embeddings are baked from a 2026 snapshot, so on a 2023--25 window they encode slow-moving corporate structure from the model's future); read it as indicative, not deployable.}
\label{tab:bt_horizons}
\centering
\small
\begin{tabular}{llrrrr}
\toprule
\textbf{OOS Window} & \textbf{Strategy} & \textbf{Return} & \textbf{Sharpe} & \textbf{Max DD} & \textbf{IR} \\
\midrule
Pre-Training$^{\ddagger}$    & Buy \& Hold      & $+33.59\%$ & $1.92$ & $-7.17\%$  & --- \\
(20 mo, 2023-07$\to$25-03)   & Top-K=10         & $+15.62\%$ & $0.70$ & $-13.47\%$ & $-0.60$ \\
\midrule
Post-Training (cleanest)     & Buy \& Hold      & $+22.18\%$ & $5.22$ & $-2.41\%$  & --- \\
(2 mo, 2026-03$\to$05)       & Top-K=10         & $\mathbf{+47.28\%}$ & $\mathbf{6.66}$ & $-2.40\%$ & $\mathbf{+4.37}$ \\
\bottomrule
\end{tabular}
\end{table}

\noindent \textbf{Headline number.} On the cleanest window --- the strictly-post-training stretch ($2026\text{-}03\text{-}20 \to 2026\text{-}05\text{-}15$), where there is neither memorisation nor meaningful survivorship leak --- the full-universe top-K=10 book returns \textbf{$+47.28\%$ vs Buy \& Hold $+22.18\%$, Sharpe $6.66$, Max DD $-2.40\%$, Information Ratio $+4.37$}. The IR of $+4.37$ is the key statistic: it says the strategy's excess return over the benchmark is large relative to its tracking error, i.e.\ the out-performance is consistent, not a single lucky bet. On the long pre-training window the same book \emph{under}-performs Buy \& Hold (IR $-0.60$); we show this rather than hide it, and attribute it partly to the survivorship/embedding caveats above and partly to genuinely different market conditions in 2023--24.

\subsection{Sector-Conditional Behaviour}

The top-K=10 strategy was further tested on five sector-specific 50-ticker subsets over the same 2-month unseen-future window. Performance is materially sector-dependent (Table~\ref{tab:bt_subsets}): the model excels on Tech and Mega-Cap (the regimes that dominate its training distribution) and \emph{loses to Buy \& Hold} on Healthcare and Financials. We report this rather than hide it: the KG embeddings encode structural relations (sector, industry, country, acquisitions) that align with the dynamics of Tech megacaps but underspecify the idiosyncratic regulatory and rate-driven dynamics of Healthcare and Financials.

\begin{table}[htbp]
\caption{Sector-conditional performance (unseen future, $2026\text{-}03\text{-}20 \to 2026\text{-}05\text{-}15$, 50-ticker subsets, top-K=10).}
\label{tab:bt_subsets}
\centering
\small
\begin{tabular}{lrrr}
\toprule
\textbf{Subset} & \textbf{Buy \& Hold} & \textbf{High-Conf top-K=10} & $\boldsymbol{\Delta}$ \\
\midrule
Mega-Cap Titans       & +14.67\% & \textbf{+19.98\%}  & $+5.3$ pp  \\
Technology Sector     & +31.36\% & \textbf{+54.37\%}  & $+23.0$ pp \\
Consumer Services     & +1.43\%  & \textbf{+7.88\%}   & $+6.5$ pp  \\
Healthcare Pioneers   & +5.55\%  & $-2.37\%$          & $-7.9$ pp  \\
Financial Giants      & +9.51\%  & $-6.28\%$          & $-15.8$ pp \\
\bottomrule
\end{tabular}
\end{table}

\noindent The combined narrative: the model carries genuine signal (CV IC $+0.1053$, $5/5$ folds positive, paired bootstrap CI excludes zero on $8/13$ Combined-minus-Tabular pairings); on the cleanest out-of-sample window the full-universe top-K=10 book converts that signal to $+47.28\%$ vs $+22.18\%$ Buy \& Hold at an Information Ratio of $+4.37$; and the signal is \emph{not uniform across sectors} (Tech/Mega-Cap win, Healthcare/Financial lose). The concentrated $K=10$ book also carries higher drawdown than a diversified index ($-2.4\%$ on the clean post-training window, $-13.5\%$ on the long pre-training window), the structural cost of trading ten names.

\subsection{Breadth Experiment: Why We Keep $K{=}10$}
A natural objection to a ten-name book is low \emph{breadth}: Grinold's Fundamental Law of Active Management states $\text{IR} \approx \text{IC}\,\sqrt{\text{breadth}}$, so more independent bets should raise the Information Ratio for the same skill. We tested this directly, widening to $K{=}20$ and adding a per-sector count cap (at most $\lceil 0.30K \rceil$ names per sector) to force the book out of its semiconductor cluster. On the clean post-training window this was \emph{worse} on every primary metric: return $+40.12\%$ vs $+47.28\%$, Sharpe $6.51$ vs $6.66$, IR $+3.32$ vs $+4.37$; and on the pre-training window dramatically worse (IR $-1.67$ vs $-0.60$). The reason is that Grinold's law assumes \emph{uniform} per-bet skill, which the sector-conditional results above falsify: forcing diversification spends capital on the model's negative-IC sectors (Healthcare, Financials), so the added "breadth" carries negative skill and drags portfolio IC down faster than $\sqrt{\text{breadth}}$ lifts consistency. We therefore retain $K{=}10$ with no sector cap; the cap remains in the codebase as a disabled, tested knob for a future, more sector-uniform model.

\section{Conclusion}
This work makes two claims that we believe matter beyond the immediate academic context. First, knowledge-graph embeddings deserve targets that match their inductive bias: a graph encoding sector, industry, country and corporate-action structure adds essentially zero predictive value at the 7-day binary direction horizon, but lifts the cross-sectional 30-day rank IC from $0.04$ (tabular-only) to $0.09$ when fused with technical and macroeconomic features under walk-forward CV. The structural signal is real once the question matches the signal. Second, the gap between a backtested model and a production trading agent is closed not by yet another classifier, but by the boring operational plumbing that makes trades reproducible, observable, and stoppable. The reproducibility hash, per-trade structured log, peak-to-trough circuit breaker, and three-bucket attribution we ship alongside the model are the difference between a research artefact and infrastructure that can carry capital.

\begin{thebibliography}{1}
\bibitem{sun2019rotate}
Z.~Sun, Z.-H. Deng, J.-Y. Nie, and J.~Tang, ``Rotate: Knowledge graph embedding by relation rotation in complex space,'' in \emph{International Conference on Learning Representations (ICLR)}, 2019.
\end{thebibliography}

\clearpage
\appendices

\section{Quantitative and Financial Glossaries}

\subsection{Systemic Risk and Kalman Beta ($\beta$)}
The Beta ($\beta$) of an asset measures its volatility in relation to the broader market index (e.g., the S\&P 500). Mathematically, it represents the slope of the security's returns relative to market returns:
\begin{equation}
\text{Returns}_{\text{asset}, t} = \alpha + \beta \cdot \text{Returns}_{\text{market}, t} + \epsilon_t
\end{equation}
A $\beta > 1.0$ indicates that the stock is highly sensitive to market movements (amplifying both gains and losses), while a $\beta < 1.0$ represents defensive, lower-sensitivity dynamics. By implementing a state-space Kalman Filter, the system estimates the time-varying, dynamic beta ($\beta_{t|t}$) recursively rather than relying on historical rolling windows which suffer from lookback lag. This enables the bot to adjust short hedges instantly as soon as a market shock occurs.

\subsection{Annualized Sharpe Ratio}
The Sharpe ratio measures the excess return earned per unit of volatility in an investment asset. It is defined as:
\begin{equation}
\text{Sharpe} = \frac{\mathbb{E}[R_p - R_f]}{\sigma_p}
\end{equation}
where $R_p$ is the portfolio return, $R_f$ is the risk-free rate, and $\sigma_p$ is the standard deviation of the portfolio returns. In our weekly backtesting framework, the Sharpe ratio is annualized by multiplying the weekly returns' average-to-standard-deviation ratio by $\sqrt{52}$:
\begin{equation}
\text{Annualized Sharpe} = \sqrt{52} \times \frac{\mu_{\text{weekly returns}}}{\sigma_{\text{weekly returns}}}
\end{equation}
A Sharpe ratio above $1.0$ is considered good, above $2.0$ is excellent, and above $3.0$ is exceptional, indicating highly stable, risk-managed returns.

\subsection{Maximum Drawdown (Max DD)}
Maximum Drawdown measures the largest peak-to-trough drop in a portfolio's capital curve before a new peak is achieved. It represents the worst-case loss scenario for an investor buying at the absolute market peak:
\begin{align}
\text{Drawdown}_t &= \frac{\text{Equity}_t - \max_{s \le t} \text{Equity}_s}{\max_{s \le t} \text{Equity}_s} \\
\text{Max DD} &= \min_{t} \text{Drawdown}_t
\end{align}
Preserving capital by compressing drawdowns is the single most important objective of our HMM + Kalman risk layer, as deep drawdowns destroy the mathematical compounding capacity of a portfolio.

\subsection{Gaussian Hidden Markov Models (HMM)}
A Hidden Markov Model assumes that the observed data sequence (in our case, daily S\&P 500 log returns) is generated by a series of hidden, unobserved states that follow a Markov chain. The state transitions are governed by a transition matrix $\mathbf{A} \in \mathbb{R}^{K \times K}$, where $a_{ij} = P(S_{t+1} = j | S_t = i)$. The emissions are modeled as continuous Gaussian distributions:
\begin{equation}
b_j(x_t) = \frac{1}{\sqrt{2\pi \sigma_j^2}} \exp\left( -\frac{(x_t - \mu_j)^2}{2\sigma_j^2} \right)
\end{equation}
where $\mu_j$ and $\sigma_j^2$ represent the mean and variance of daily index log returns in state $j$. By training the model, the system decodes the market state dynamically, shutting down short positions in the low-volatility bull regime to protect the portfolio from short squeeze drag.

\subsection{Transaction Friction and Slippage}
In theoretical backtests, assets are assumed to trade at the exact closing price. In real markets, execution incurs three direct frictions:
\begin{enumerate}
    \item \textbf{Brokerage Commissions}: The flat fee charged by the execution venue.
    \item \textbf{Bid-Ask Spread}: The difference between the highest price a buyer is willing to pay and the lowest price a seller is willing to accept.
    \item \textbf{Market Slippage}: The price impact caused by executing large volumes that consume available order book liquidity.
\end{enumerate}
By charging a strict 10 bps fee rate on every rebalancing turnover delta, our backtests represent a highly realistic, institutional-grade simulation. The differential optimizer reduces this friction directly by trading only the required weight deltas rather than executing full liquidations, saving substantial transaction drag.

\end{document}
